% ============================================
% ALGEBRA HOMEWORK TEMPLATE
% Clean, minimalist design for math assignments
% Compatible with pdfLaTeX and XeLaTeX
% ============================================

\documentclass[12pt,a4paper]{article}

% ---------- Encoding and Language ----------
\usepackage[utf8]{inputenc}          % UTF-8 encoding (use utf8x if extended Cyrillic needed)
\usepackage[T2A]{fontenc}           % Cyrillic font encoding (for pdfLaTeX)
\usepackage[russian]{babel}         % Russian language support

% ---------- Page Geometry ----------
\usepackage[
  left=2cm,
  right=2cm,
  top=2cm,
  bottom=2cm,
  headsep=0.5cm,
  footskip=1cm
]{geometry}

% ---------- Mathematics ----------
\usepackage{amsmath}                % Core math features
\usepackage{amssymb}                % Additional math symbols
\usepackage{amsthm}                 % Theorem-like environments

% ---------- Lists (Compact & Multi-level) ----------
\usepackage{enumitem}
% Set compact defaults for all lists
\setlist{
  noitemsep,                        % No extra space between items
  topsep=0.5ex,                     % Minimal space above/below list
  partopsep=0pt,                    % No extra space when starting a paragraph
  parsep=0pt,                       % No space between paragraphs within items
  leftmargin=2em,                   % Consistent indentation
  labelsep=0.5em                    % Space between label and text
}
% Fine-tune specific list types
\setlist[itemize]{label=---}        % Em-dash bullets
\setlist[enumerate,1]{label=\arabic*.} % First level: 1.
\setlist[enumerate,2]{label=\alph*)}  % Second level: a)
\setlist[enumerate,3]{label=\roman*.} % Third level: i.
\setlist[description]{font=\normalfont\itshape} % Description labels in italics

% ---------- Problem Environment ----------
\newcounter{task}                   % Custom counter for problems
\newcommand{\taskname}{}            % Placeholder for task label

% Define the task environment
\newenvironment{task}[1][]{%
  \refstepcounter{task}%            % Increment and reference the counter
  \par\vspace{0.3\baselineskip}%    % Tiny space before problem
  \noindent\textbf{\thetask.}%      % Bold number with period
  \if\relax\detokenize{#1}\relax   % Check if optional title is provided
    \ \ignorespaces                 % Just space after number
  \else
    \ \textbf{#1}.\ \ignorespaces   % Space, bold title, period, space
  \fi
}{%
  \par\vspace{0.2\baselineskip}%    % Tiny space after problem
}

% ---------- Header Configuration ----------
\makeatletter
\renewcommand{\@maketitle}{%
  % Top row: Sheet number (left) and Student name (right)
  \begin{flushleft}
    \textbf{Листок \#\sheetnumber}%
    \hfill%
    \textbf{\studentname}%
  \end{flushleft}
  \vspace{-0.8\baselineskip}%      % Tighten vertical spacing
  
  % Bottom row: Assignment title
  \begin{flushleft}
    \textit{\assignmenttitle}%
  \end{flushleft}
  \vspace{-0.5\baselineskip}%      % Additional tightening
}
\makeatother

% Define header variables
\newcommand{\sheetnumber}[1]{\def\sheetnumber{#1}}
\newcommand{\studentname}[1]{\def\studentname{#1}}
\newcommand{\assignmenttitle}[1]{\def\assignmenttitle{#1}}

% ---------- Other Settings ----------
\date{}                             % Suppress automatic date
\pagestyle{plain}                   % Simple page numbering
\setlength{\parindent}{0pt}         % No paragraph indentation (modern style)
\setlength{\parskip}{0.5ex}         % Small paragraph spacing

% ============================================
% DOCUMENT BEGIN
% ============================================

\begin{document}

% ---------- Custom Header ----------
% Set your values here:
\sheetnumber{5}                     % Replace with actual sheet number
\studentname{Имя Фамилия}           % Replace with student name
\assignmenttitle{Алгебраические структуры и группы} % Replace with assignment title
\maketitle

% ---------- Example Problems ----------

% --- Problem 1 ---
\begin{task}
  Докажите, что множество $\mathbb{Z}$ с операцией сложения образует абелеву группу.
\end{task}

% --- Problem 2 ---
\begin{task}
  Пусть $G$ — группа. Докажите, что для любых $a,b \in G$ выполняется:
  \[
    (ab)^{-1} = b^{-1}a^{-1}.
  \]
\end{task}

% --- Problem 3 with enumerated sub-items ---
\begin{task}
  Определите, являются ли следующие множества группами относительно указанной операции:
  \begin{enumerate}
    \item Множество $\mathbb{R}^+$ с операцией умножения.
    \item Множество $\mathbb{R}$ с операцией $a * b = a + b + ab$.
    \item Множество матриц вида $\begin{pmatrix} a & b \\ 0 & 1 \end{pmatrix}$ ($a \neq 0$) с матричным умножением.
  \end{enumerate}
\end{task}

% --- Problem 4 with itemize ---
\begin{task}
  Приведите примеры следующих алгебраических структур:
  \begin{itemize}
    \item Полугруппы, не являющейся моноидом.
    \item Неабелевой группы порядка 6.
    \item Кольца с делителями нуля.
  \end{itemize}
\end{task}

% --- Problem 5 with description list ---
\begin{task}
  Сопоставьте каждую структуру с её определением:
  \begin{description}
    \item[Группа] Множество с ассоциативной бинарной операцией, в котором существует нейтральный элемент и для каждого элемента существует обратный.
    \item[Кольцо] Множество с двумя бинарными операциями, образующее абелеву группу по сложению и моноид по умножению, связанных дистрибутивным законом.
    \item[Поле] Коммутативное кольцо с единицей, в котором каждый ненулевой элемент обратим.
  \end{description}
\end{task}

% --- Problem 6 with nested enumeration ---
\begin{task}
  Рассмотрите группу перестановок $S_3$:
  \begin{enumerate}
    \item Найдите все подгруппы группы $S_3$:
    \begin{enumerate}
      \item определите порядок каждой подгруппы;
      \item проверьте, какие из них являются нормальными.
    \end{enumerate}
    \item Постройте таблицу Кэли для факторгруппы, если она существует.
  \end{enumerate}
\end{task}

% --- Problem 7 with label/reference ---
\begin{task}
  \label{task:lagrange}
  Сформулируйте и докажите теорему Лагранжа о порядке подгруппы конечной группы.
\end{task}

% Referencing example (uncomment to use):
% Как показано в задаче \ref{task:lagrange}, порядок подгруппы делит порядок группы.

\end{document}
